\section{Conclusion}
\label{sec:conclusion}

We measured the residual effects of four common categories of system prompt on five unrelated benchmarks and three open instruction-tuned models, using paired comparisons against a minimal baseline, length-matched controls, two scorers and a full sampling replicate. Safety prompts do not reduce reasoning accuracy; they add hedging and modestly improve truthfulness. Persona prompts double output length and cost about 3\pp of accuracy. Format prompts make answers parseable ($+19.6$\pp strict) but not more correct ($-2.5$\pp lenient), and rigid schemas can shorten reasoning. Constraint prompts are consistently the most harmful ($-9.0$\pp), and catastrophically so ($-61$\pp) when they forbid visible reasoning. The key takeaway is that, when auditing a system prompt, constraints on length and style matter far more than safety boilerplate.

Several directions remain open. First, frontier API models may differ, especially in refusal behaviour. Second, varying the {\em strength} of each category would give dose--response curves rather than category averages. Third, two-stage ``reason freely, then format'' prompts could keep the parseability benefit of format prompts without their reasoning cost. Fourth, an LLM judge could score prose answers and remove the dependence on answer extraction. Finally, the scope refusals we observed from a neutral deployment blurb deserve a systematic study.
